% Options for packages loaded elsewhere
\PassOptionsToPackage{unicode}{hyperref}
\PassOptionsToPackage{hyphens}{url}
\documentclass[
]{article}
\usepackage{xcolor}
\usepackage{amsmath,amssymb}
\setcounter{secnumdepth}{-\maxdimen} % remove section numbering
\usepackage{iftex}
\ifPDFTeX
  \usepackage[T1]{fontenc}
  \usepackage[utf8]{inputenc}
  \usepackage{textcomp} % provide euro and other symbols
\else % if luatex or xetex
  \usepackage{unicode-math} % this also loads fontspec
  \defaultfontfeatures{Scale=MatchLowercase}
  \defaultfontfeatures[\rmfamily]{Ligatures=TeX,Scale=1}
\fi
\usepackage{lmodern}
\ifPDFTeX\else
  % xetex/luatex font selection
\fi
% Use upquote if available, for straight quotes in verbatim environments
\IfFileExists{upquote.sty}{\usepackage{upquote}}{}
\IfFileExists{microtype.sty}{% use microtype if available
  \usepackage[]{microtype}
  \UseMicrotypeSet[protrusion]{basicmath} % disable protrusion for tt fonts
}{}
\makeatletter
\@ifundefined{KOMAClassName}{% if non-KOMA class
  \IfFileExists{parskip.sty}{%
    \usepackage{parskip}
  }{% else
    \setlength{\parindent}{0pt}
    \setlength{\parskip}{6pt plus 2pt minus 1pt}}
}{% if KOMA class
  \KOMAoptions{parskip=half}}
\makeatother
\usepackage{color}
\usepackage{fancyvrb}
\newcommand{\VerbBar}{|}
\newcommand{\VERB}{\Verb[commandchars=\\\{\}]}
\DefineVerbatimEnvironment{Highlighting}{Verbatim}{commandchars=\\\{\}}
% Add ',fontsize=\small' for more characters per line
\newenvironment{Shaded}{}{}
\newcommand{\AlertTok}[1]{\textcolor[rgb]{1.00,0.00,0.00}{\textbf{#1}}}
\newcommand{\AnnotationTok}[1]{\textcolor[rgb]{0.38,0.63,0.69}{\textbf{\textit{#1}}}}
\newcommand{\AttributeTok}[1]{\textcolor[rgb]{0.49,0.56,0.16}{#1}}
\newcommand{\BaseNTok}[1]{\textcolor[rgb]{0.25,0.63,0.44}{#1}}
\newcommand{\BuiltInTok}[1]{\textcolor[rgb]{0.00,0.50,0.00}{#1}}
\newcommand{\CharTok}[1]{\textcolor[rgb]{0.25,0.44,0.63}{#1}}
\newcommand{\CommentTok}[1]{\textcolor[rgb]{0.38,0.63,0.69}{\textit{#1}}}
\newcommand{\CommentVarTok}[1]{\textcolor[rgb]{0.38,0.63,0.69}{\textbf{\textit{#1}}}}
\newcommand{\ConstantTok}[1]{\textcolor[rgb]{0.53,0.00,0.00}{#1}}
\newcommand{\ControlFlowTok}[1]{\textcolor[rgb]{0.00,0.44,0.13}{\textbf{#1}}}
\newcommand{\DataTypeTok}[1]{\textcolor[rgb]{0.56,0.13,0.00}{#1}}
\newcommand{\DecValTok}[1]{\textcolor[rgb]{0.25,0.63,0.44}{#1}}
\newcommand{\DocumentationTok}[1]{\textcolor[rgb]{0.73,0.13,0.13}{\textit{#1}}}
\newcommand{\ErrorTok}[1]{\textcolor[rgb]{1.00,0.00,0.00}{\textbf{#1}}}
\newcommand{\ExtensionTok}[1]{#1}
\newcommand{\FloatTok}[1]{\textcolor[rgb]{0.25,0.63,0.44}{#1}}
\newcommand{\FunctionTok}[1]{\textcolor[rgb]{0.02,0.16,0.49}{#1}}
\newcommand{\ImportTok}[1]{\textcolor[rgb]{0.00,0.50,0.00}{\textbf{#1}}}
\newcommand{\InformationTok}[1]{\textcolor[rgb]{0.38,0.63,0.69}{\textbf{\textit{#1}}}}
\newcommand{\KeywordTok}[1]{\textcolor[rgb]{0.00,0.44,0.13}{\textbf{#1}}}
\newcommand{\NormalTok}[1]{#1}
\newcommand{\OperatorTok}[1]{\textcolor[rgb]{0.40,0.40,0.40}{#1}}
\newcommand{\OtherTok}[1]{\textcolor[rgb]{0.00,0.44,0.13}{#1}}
\newcommand{\PreprocessorTok}[1]{\textcolor[rgb]{0.74,0.48,0.00}{#1}}
\newcommand{\RegionMarkerTok}[1]{#1}
\newcommand{\SpecialCharTok}[1]{\textcolor[rgb]{0.25,0.44,0.63}{#1}}
\newcommand{\SpecialStringTok}[1]{\textcolor[rgb]{0.73,0.40,0.53}{#1}}
\newcommand{\StringTok}[1]{\textcolor[rgb]{0.25,0.44,0.63}{#1}}
\newcommand{\VariableTok}[1]{\textcolor[rgb]{0.10,0.09,0.49}{#1}}
\newcommand{\VerbatimStringTok}[1]{\textcolor[rgb]{0.25,0.44,0.63}{#1}}
\newcommand{\WarningTok}[1]{\textcolor[rgb]{0.38,0.63,0.69}{\textbf{\textit{#1}}}}
\usepackage{longtable,booktabs,array}
\newcounter{none} % for unnumbered tables
\usepackage{calc} % for calculating minipage widths
% Correct order of tables after \paragraph or \subparagraph
\usepackage{etoolbox}
\makeatletter
\patchcmd\longtable{\par}{\if@noskipsec\mbox{}\fi\par}{}{}
\makeatother
% Allow footnotes in longtable head/foot
\IfFileExists{footnotehyper.sty}{\usepackage{footnotehyper}}{\usepackage{footnote}}
\makesavenoteenv{longtable}
\setlength{\emergencystretch}{3em} % prevent overfull lines
\providecommand{\tightlist}{%
  \setlength{\itemsep}{0pt}\setlength{\parskip}{0pt}}
\usepackage{bookmark}
\IfFileExists{xurl.sty}{\usepackage{xurl}}{} % add URL line breaks if available
\urlstyle{same}
\hypersetup{
  hidelinks,
  pdfcreator={LaTeX via pandoc}}

\author{}
\date{}

\begin{document}

\section{BH Thermodynamics Programme Paper 7 Supplement: Geometric
Observations on the Second-Order Correction of the α
Identity}\label{bh-thermodynamics-programme-paper-7-supplement-geometric-observations-on-the-second-order-correction-of-the-ux3b1-identity}

\textbf{Author}: Noriaki Kihara \textbf{Affiliation}: WF System Co.,
Ltd. \textbf{ORCID}:
\href{https://orcid.org/0009-0004-6753-4020}{0009-0004-6753-4020}
\textbf{Version}: v1 (draft) \textbf{Date}: May 1, 2026
\textbf{License}: CC BY 4.0

\begin{center}\rule{0.5\linewidth}{0.5pt}\end{center}

\subsection{Character of This
Document}\label{character-of-this-document}

\textbf{This is an observation paper, not a proof paper.}

For the α identity \(\alpha^{-1} = 137 + (\pi^2/2)\alpha\) in Paper 7
{[}BH7{]}, we provide one geometric interpretive framework concerning
the coefficient \(\pi^2/2\) and the residual at the level of 8.7 ppb
against observation.

This document does NOT claim:

\begin{itemize}
\tightlist
\item
  That \(\pi^2/2\) admits no interpretation other than the one presented
  here
\item
  That the residual of 8.7 ppb is fully derived under our interpretation
\item
  That other interpretive possibilities are excluded
\item
  That this document ``proves,'' ``strengthens,'' or ``extends'' the
  claims of Paper 7
\end{itemize}

What this document presents is \textbf{one reading}. The possibility
that the same equations and numbers admit other readings is not denied.

The claims of Paper 7 do not depend on this document. The numerical fact
observed in Paper 7 (that \(\alpha^{-1} = 137 + (\pi^2/2)\alpha\) holds
to 8.7 ppb precision) does not change regardless of whether one accepts
or rejects this document's interpretation.

\subsection{Abstract}\label{abstract}

In Paper 7 {[}BH7{]}, a geometric identity using unit cube-packing of a
4-dimensional ball was reported:

\[\alpha^{-1} = 137 + \frac{\pi^2}{2}\alpha,\]

agreeing with the observed value \(\alpha^{-1} = 137.035999177\) at
relative precision \(8.7 \times 10^{-8}\) (8.7 ppb). This document
offers \textbf{one geometric interpretation} of this structural
relation. Specifically:

\begin{enumerate}
\def\labelenumi{\arabic{enumi}.}
\tightlist
\item
  \(\pi^2/2\) admits a reading as the ``normalized boundary measure
  \(A_{S^3}(R)/(4R^3)\)'' of the 4D ball boundary \(S^3\), rather than
  as the ``unit ball volume \(V_4(1)\)'' (§2).
\item
  Under this reading, the volume deficit \(G(R)\) of cube-covering of a
  4D ball admits the asymptotic expansion (§3):
  \[G(R) = \frac{16\pi}{3}R^3 - \left(\frac{3\pi^2}{4} + \frac{9\pi}{2}\right)R^2 + O(R).\]
\item
  The second-coefficient is analytically \textbf{negative} in this
  expansion (§4).
\item
  Numerical experiments (R = 50--300) are consistent with this expansion
  (§5).
\item
  The back-calculated correction \(\delta \approx 0.001619\) from the
  observed residual is consistent in sign and magnitude with this
  geometric structure (§6).
\end{enumerate}

This document does not provide a complete derivation of \(\delta\). Item
5 remains an observation; the normalization rule mapping the \(R^2\)
deficit term to the coefficient term in the α equation is left as an
open question (§7).

\subsection{§1 Introduction}\label{introduction}

\subsubsection{1.1 Context of Paper 7}\label{context-of-paper-7}

Paper 7 {[}BH7{]} considers unit hypercube cells centered on integer
lattice points in the 4D Euclidean lattice, and establishes the
self-consistent identity:

\[\alpha^{-1} = 137 + \frac{\pi^2}{2}\alpha \tag{1.1}\]

based on the count \(N(1) = 137\) of cells fully contained in a 4D ball
of radius \(R = 3\), and the unit 4-ball volume \(V_4(1) = \pi^2/2\).
The positive root yields \(\alpha = 7.29735194 \times 10^{-3}\),
agreeing with the CODATA 2022 value
\(\alpha = 7.2973525693 \times 10^{-3}\) at relative precision
\(8.7 \times 10^{-8}\).

\subsubsection{1.2 The Residual Question}\label{the-residual-question}

Paper 7 left the origin of the residual \(6.3 \times 10^{-10}\) as an
open problem, identifying it as a candidate for an \(\alpha^3\)-order
geometric correction with coefficient
\(c_3 \approx 1.6 \times 10^{-3}\).

\subsubsection{1.3 Anticipated Concerns}\label{anticipated-concerns}

The structure of Paper 7 may be subject to such concerns as:

\begin{itemize}
\tightlist
\item
  \textbf{Concern A}: ``\(\pi^2/2\) as the unit ball volume may have
  been chosen post-hoc to match observation.''
\item
  \textbf{Concern B}: ``Without a geometric origin for the residual, the
  coefficient \(\pi^2/2\) cannot be excluded as numerical fitting.''
\end{itemize}

This document offers a \textbf{partial response} to these concerns. It
is not a complete refutation but an observation that one geometric
interpretation is consistent with the structure of Paper 7.

\subsubsection{1.4 Document Structure}\label{document-structure}

§2 presents an alternative reading of \(\pi^2/2\) (as the boundary
measure normalization core). §3 derives the asymptotic expansion of the
volume deficit \(G(R)\). §4 shows that the second coefficient is
analytically negative. §5 confirms consistency with numerical
experiments. §6 records the relation to the observed residual as
observation. §7 explicitly acknowledges the limitations and open
questions. §8 organizes implications for Paper 7.

\begin{center}\rule{0.5\linewidth}{0.5pt}\end{center}

\subsection{\texorpdfstring{§2 An Alternative Geometric Reading of
\(\pi^2/2\)}{§2 An Alternative Geometric Reading of \textbackslash pi\^{}2/2}}\label{an-alternative-geometric-reading-of-pi22}

\subsubsection{2.1 Numerical Equivalence of Two
Interpretations}\label{numerical-equivalence-of-two-interpretations}

In Paper 7, \(\pi^2/2\) was introduced as the volume of the unit 4-ball:

\[V_4(R) = \frac{\pi^2}{2}R^4, \quad V_4(1) = \frac{\pi^2}{2}. \tag{2.1}\]

However, \(\pi^2/2\) also appears naturally from a different route, via
the boundary measure of the 4-ball:

\[A_{S^3}(R) = \frac{d}{dR}V_4(R) = 2\pi^2 R^3. \tag{2.2}\]

Normalizing this by \(4R^3\):

\[\frac{A_{S^3}(R)}{4R^3} = \frac{2\pi^2 R^3}{4R^3} = \frac{\pi^2}{2}. \tag{2.3}\]

Numerically, \(V_4(1) = A_{S^3}(R)/(4R^3) = \pi^2/2\); the two are
equal.

\subsubsection{2.2 The Two Readings Differ in
Meaning}\label{the-two-readings-differ-in-meaning}

Although the numerical value is the same, the two readings differ in
their \textbf{physical/geometric meaning}:

{\def\LTcaptype{none} % do not increment counter
\begin{longtable}[]{@{}
  >{\raggedright\arraybackslash}p{(\linewidth - 4\tabcolsep) * \real{0.3333}}
  >{\raggedright\arraybackslash}p{(\linewidth - 4\tabcolsep) * \real{0.3333}}
  >{\raggedright\arraybackslash}p{(\linewidth - 4\tabcolsep) * \real{0.3333}}@{}}
\toprule\noalign{}
\begin{minipage}[b]{\linewidth}\raggedright
Reading
\end{minipage} & \begin{minipage}[b]{\linewidth}\raggedright
Meaning
\end{minipage} & \begin{minipage}[b]{\linewidth}\raggedright
Scale dependence
\end{minipage} \\
\midrule\noalign{}
\endhead
\bottomrule\noalign{}
\endlastfoot
\(V_4(1)\) & Volume of the unit 4-ball & Choice of ``unit'' present \\
\(A_{S^3}(R)/(4R^3)\) & Boundary measure of 4-ball normalized by \(R^3\)
& \textbf{Independent of R} (same value for any R) \\
\end{longtable}
}

The second reading does not specify any particular \(R\). It yields the
same value \(\pi^2/2\) for any \(R\). This carries the property of a
\textbf{scale-invariant geometric core}.

\subsubsection{2.3 The Reading Proposed by This
Document}\label{the-reading-proposed-by-this-document}

This document proposes that \(\pi^2/2\) in Paper 7 (1.1) admits the
reading

\begin{quote}
\textbf{as the ``normalization core of the boundary measure
\(A_{S^3}(R)/(4R^3)\) of the 4D ball'' rather than as the ``volume
\(V_4(1)\) of the unit 4D ball,''}
\end{quote}

as \textbf{one possible interpretation}.

Caveats:

\begin{itemize}
\tightlist
\item
  This is not the ``more correct'' interpretation (since the numerical
  values are identical, mathematics cannot select between them).
\item
  It simply provides the observation that \(\pi^2/2\) appears as a
  scale-invariant quantity, free of any arbitrary unit choice.
\item
  Re-reading the structure of Paper 7 under this lens leads naturally to
  the developments of §3 onward.
\end{itemize}

\begin{center}\rule{0.5\linewidth}{0.5pt}\end{center}

\subsection{\texorpdfstring{§3 Asymptotic Expansion of
\(G(R)\)}{§3 Asymptotic Expansion of G(R)}}\label{asymptotic-expansion-of-gr}

To analyze the structure of Paper 7 under this interpretive framework,
we quantify the geometric mismatch between the 4D ball and the cubic
lattice.

\subsubsection{3.1 Definition of Volume
Deficit}\label{definition-of-volume-deficit}

On the 4D Euclidean integer lattice \(\mathbb{Z}^4\), place a unit
hypercube cell of side 1 centered at each lattice point. Let
\(N_{\text{full}}(R)\) denote the number of cells fully contained in the
4D ball \(B_R\) of radius \(R\).

Define the volume deficit \(G(R)\) as:

\[G(R) := V_4(R) - N_{\text{full}}(R) = \frac{\pi^2}{2}R^4 - N_{\text{full}}(R). \tag{3.1}\]

This is ``the residual when ball volume minus fully-contained cell
volume is taken,'' corresponding to the contribution of cells partially
cut by the spherical surface.

\subsubsection{3.2 Hypercube Support
Function}\label{hypercube-support-function}

The outward support function of a unit hypercube with respect to the
unit normal \(u = (u_1, u_2, u_3, u_4) \in S^3\) on the spherical
surface is defined by:

\[\tau(u) = \frac{1}{2}\sum_{i=1}^{4} |u_i|. \tag{3.2}\]

This represents the ``outward protrusion of the hypercube'' in direction
\(u\).

\subsubsection{3.3 Sphere Integrals}\label{sphere-integrals}

To compute the asymptotic expansion coefficients, we introduce two basic
integrals on \(S^3\).

\paragraph{\texorpdfstring{Lemma 3.1:
\(\int_{S^3} \tau(u) \, d\Omega = 16\pi/3\)}{Lemma 3.1: \textbackslash int\_\{S\^{}3\} \textbackslash tau(u) \textbackslash, d\textbackslash Omega = 16\textbackslash pi/3}}\label{lemma-3.1-int_s3-tauu-domega-16pi3}

\textbf{Proof}: From the symmetry of
\(\tau(u) = (1/2)\sum_{i=1}^4 |u_i|\),

\[\int_{S^3} \tau(u) \, d\Omega = 2 \int_{S^3} |u_1| \, d\Omega.\]

Parameterize \(S^3\) with \(u_1 = \cos\theta\), \(\theta \in [0, \pi]\).
The volume element of \(S^3\) is
\(\sin^2\theta \, d\theta \, d\Omega_{S^2}\), and the surface area of
\(S^2\) is \(4\pi\). Thus:

\[\int_{S^3} |u_1| \, d\Omega = 4\pi \int_0^\pi |\cos\theta| \sin^2\theta \, d\theta = 8\pi \int_0^{\pi/2} \cos\theta \sin^2\theta \, d\theta = \frac{8\pi}{3}.\]

Therefore \(\int_{S^3} \tau \, d\Omega = 2 \cdot 8\pi/3 = 16\pi/3\).
\(\square\)

\paragraph{\texorpdfstring{Lemma 3.2:
\(\int_{S^3} \tau(u)^2 \, d\Omega = \pi^2/2 + 3\pi\)}{Lemma 3.2: \textbackslash int\_\{S\^{}3\} \textbackslash tau(u)\^{}2 \textbackslash, d\textbackslash Omega = \textbackslash pi\^{}2/2 + 3\textbackslash pi}}\label{lemma-3.2-int_s3-tauu2-domega-pi22-3pi}

\textbf{Proof}:
\(\tau^2 = (1/4)(\sum |u_i|)^2 = (1/4)[\sum u_i^2 + 2\sum_{i<j} |u_i||u_j|]\).

First term: \(\sum_{i=1}^4 u_i^2 = 1\) (definition of unit sphere).
Surface area of \(S^3\) is \(2\pi^2\), so

\[\int_{S^3} \sum_i u_i^2 \, d\Omega = 2\pi^2.\]

Second term: by symmetry, all 6 pairs \(i \neq j\) have equal
\(\int |u_i||u_j| \, d\Omega\). Parameterize \(S^3\) with
\(u_1 = \cos\theta_1\), \(u_2 = \sin\theta_1 \cos\theta_2\),
\(u_3 = \sin\theta_1 \sin\theta_2 \cos\theta_3\),
\(u_4 = \sin\theta_1 \sin\theta_2 \sin\theta_3\) (volume element
\(\sin^2\theta_1 \sin\theta_2 \, d\theta_1 d\theta_2 d\theta_3\)):

\[\int_{S^3} |u_1||u_2| \, d\Omega = 2\pi \cdot \int_0^\pi |\cos\theta_1| \sin^3\theta_1 \, d\theta_1 \cdot \int_0^\pi |\cos\theta_2| \sin\theta_2 \, d\theta_2.\]

\(\int_0^\pi |\cos\theta_1| \sin^3\theta_1 \, d\theta_1 = 2 \int_0^{\pi/2} \cos\theta_1 \sin^3\theta_1 \, d\theta_1 = 1/2\).

\(\int_0^\pi |\cos\theta_2| \sin\theta_2 \, d\theta_2 = 2 \int_0^{\pi/2} \cos\theta_2 \sin\theta_2 \, d\theta_2 = 1\).

So \(\int_{S^3} |u_1||u_2| \, d\Omega = 2\pi \cdot 1/2 \cdot 1 = \pi\).

Therefore \(\sum_{i<j} \int |u_i||u_j| \, d\Omega = 6\pi\), and

\[\int_{S^3} \tau^2 \, d\Omega = \frac{1}{4}[2\pi^2 + 2 \cdot 6\pi] = \frac{\pi^2}{2} + 3\pi. \quad \square\]

\subsubsection{\texorpdfstring{3.4 Asymptotic Expansion of
\(G(R)\)}{3.4 Asymptotic Expansion of G(R)}}\label{asymptotic-expansion-of-gr-1}

We adopt the approximation that, on the spherical boundary
\(\partial B_R\) in direction \(u\), the hypercube cell is shaved inward
by \(\tau(u)\). This makes the effective radius in direction \(u\) equal
to \(R - \tau(u)\).

\[G(R) \approx \frac{1}{4}\int_{S^3}\left[R^4 - (R - \tau(u))^4\right] d\Omega. \tag{3.3}\]

Here the factor \(1/4\) is the normalization derived from
\(V_4(R) = \pi^2 R^4/2\) and \(A_{S^3}(R) = 2\pi^2 R^3\) (the relation
\(dV/dR = A_{S^3}/4 \cdot\) spherical mean).

Binomial expansion:

\[R^4 - (R - \tau)^4 = 4R^3 \tau - 6R^2 \tau^2 + 4R \tau^3 - \tau^4.\]

Taking \(\int_{S^3}\) and multiplying by \(1/4\):

\[G(R) = R^3 \int_{S^3}\tau \, d\Omega - \frac{3}{2}R^2 \int_{S^3}\tau^2 \, d\Omega + R \int_{S^3}\tau^3 \, d\Omega - \frac{1}{4}\int_{S^3}\tau^4 \, d\Omega.\]

Substituting Lemmas 3.1 and 3.2:

\[\boxed{G(R) = \frac{16\pi}{3}R^3 - \left(\frac{3\pi^2}{4} + \frac{9\pi}{2}\right)R^2 + O(R).} \tag{3.4}\]

\subsubsection{3.5 Meaning of the Main
Term}\label{meaning-of-the-main-term}

The first term \((16\pi/3)R^3\) is the surface area of the 4-ball
boundary \(S^3\), which is \(2\pi^2 R^3\), multiplied by the average
boundary thickness of the hypercube. This admits the standard
interpretation as ``the main contribution from the continuous ball's
boundary layer.''

\begin{center}\rule{0.5\linewidth}{0.5pt}\end{center}

\subsection{§4 Observation on the Sign of the Second
Term}\label{observation-on-the-sign-of-the-second-term}

\subsubsection{4.1 Sign Analysis}\label{sign-analysis}

The coefficient of the second term in (3.4) is:

\[-\left(\frac{3\pi^2}{4} + \frac{9\pi}{2}\right) \approx -7.402 - 14.137 = -21.539. \tag{4.1}\]

This is a negative value. In general,
\(\int_{S^3}\tau^2 \, d\Omega > 0\), and since the coefficient \(-3/2\)
stands in front, \textbf{this term is necessarily negative under this
interpretive framework}.

\subsubsection{4.2 Interpretation}\label{interpretation}

Adopting the approximation model ``the spherical boundary is shaved
inward by the hypercube structure'':

\begin{itemize}
\tightlist
\item
  First term (\(R^3\)): main contribution from the continuous ball's
  boundary
\item
  Second term (\(R^2\)): second-order effect of the ``shaving''
\end{itemize}

The negative sign of this second term means, \textbf{under this
approximation model}, that estimating the volume deficit using only the
main term overestimates relative to the continuous ball approximation.

Caveats:

\begin{itemize}
\tightlist
\item
  This may be read as ``a fundamental property of the mismatch between
  ball and hypercube''
\item
  However, this does not claim ``the \(\pi^2/2\) in Paper 7 admits no
  explanation other than this interpretation''
\item
  The possibility that similar asymptotic expansions can be obtained
  from other approximation models or other geometric constructions is
  not excluded.
\end{itemize}

\begin{center}\rule{0.5\linewidth}{0.5pt}\end{center}

\subsection{§5 Numerical Verification}\label{numerical-verification}

\subsubsection{5.1 Numerical Experiment
Setup}\label{numerical-experiment-setup}

For integer \(R \in \{50, 100, 200, 300\}\), compute:

\begin{enumerate}
\def\labelenumi{\arabic{enumi}.}
\tightlist
\item
  \(N_{\text{full}}(R)\): number of integer centers
  \(n \in \mathbb{Z}^4\) satisfying
  \(\sum_{i=1}^4 (n_i + \epsilon_i)^2 \le R^2 \, (\forall \epsilon_i \in \{\pm 1/2\})\).
\item
  \(V_4(R) = \pi^2 R^4/2\).
\item
  \(G(R) = V_4(R) - N_{\text{full}}(R)\).
\item
  \(G(R)/R^3\) (value normalized to the main term).
\item
  \([G(R) - (16\pi/3)R^3]/R^2\) (second-coefficient after removing main
  term).
\end{enumerate}

The implementation is in Python (numpy + integer enumeration). Code is
shown in Appendix A.

\subsubsection{5.2 Results}\label{results}

{\def\LTcaptype{none} % do not increment counter
\begin{longtable}[]{@{}
  >{\raggedleft\arraybackslash}p{(\linewidth - 6\tabcolsep) * \real{0.2500}}
  >{\raggedleft\arraybackslash}p{(\linewidth - 6\tabcolsep) * \real{0.2500}}
  >{\raggedleft\arraybackslash}p{(\linewidth - 6\tabcolsep) * \real{0.2500}}
  >{\raggedleft\arraybackslash}p{(\linewidth - 6\tabcolsep) * \real{0.2500}}@{}}
\toprule\noalign{}
\begin{minipage}[b]{\linewidth}\raggedleft
\(R\)
\end{minipage} & \begin{minipage}[b]{\linewidth}\raggedleft
Filling rate \(N_{\text{full}}/V_4\)
\end{minipage} & \begin{minipage}[b]{\linewidth}\raggedleft
\(G(R)/R^3\)
\end{minipage} & \begin{minipage}[b]{\linewidth}\raggedleft
\([G(R) - (16\pi/3)R^3]/R^2\)
\end{minipage} \\
\midrule\noalign{}
\endhead
\bottomrule\noalign{}
\endlastfoot
50 & 0.9339 & 16.3130 & \(-22.11\) \\
100 & 0.9665 & 16.5529 & \(-20.22\) \\
200 & 0.9831 & 16.6595 & \(-19.13\) \\
300 & 0.9887 & 16.6838 & \(-21.42\) \\
\end{longtable}
}

Theoretical values: - \(G(R)/R^3 \to 16\pi/3 \approx 16.755\) - Second
coefficient \(\to -(3\pi^2/4 + 9\pi/2) \approx -21.539\)

\subsubsection{5.3 Observation}\label{observation}

\begin{itemize}
\tightlist
\item
  \(G(R)/R^3\) asymptotically approaches \(16\pi/3\) as \(R\) increases.
\item
  The second coefficient fluctuates in the range \(-20\) to \(-22\),
  distributed around the theoretical value \(-21.539\).
\item
  The fluctuation amplitude is consistent with finite-size effects of
  the discrete lattice.
\end{itemize}

The numerical experiment is consistent with the asymptotic expansion
(3.4). This is an observed fact. However, consistency does not show
``this interpretation is correct''; it only shows ``this interpretation
is not numerically inconsistent.''

\begin{center}\rule{0.5\linewidth}{0.5pt}\end{center}

\subsection{§6 Relation to the Observed
Residual}\label{relation-to-the-observed-residual}

\subsubsection{6.1 Back-Calculation from Observed
Value}\label{back-calculation-from-observed-value}

Introduce a correction term \(\delta\) in the equation (1.1) of Paper 7:

\[\alpha^{-1} = 137 + \left(\frac{\pi^2}{2} - \delta\right)\alpha. \tag{6.1}\]

Back-calculating from CODATA 2022 value \(\alpha^{-1} = 137.035999177\):

\[\delta = \frac{\pi^2}{2} - \alpha^{-1}(\alpha^{-1} - 137) \approx 0.001619. \tag{6.2}\]

\subsubsection{6.2 Consistency in Magnitude and
Sign}\label{consistency-in-magnitude-and-sign}

The second term in (3.4) is an \(R^2\)-order geometric correction. The
\(\delta\) in (6.2), on the other hand, is a dimensionless coefficient
that multiplies the coefficient term in the α equation directly.

To establish a direct correspondence between the two, a normalization
rule that converts the \(R^2\) deficit term into ``the coefficient term
in the α equation'' is needed (discussed in §7). This rule is not
formulated in this document.

However, the following observations can be recorded:

\begin{itemize}
\tightlist
\item
  The coefficient \(-21.539\) of the second term in (3.4) is
  \textbf{negative}.
\item
  The \(\delta\) in (6.2) is also a quantity \textbf{subtracted} from
  \(\pi^2/2\).
\item
  Their signs are consistent.
\item
  In magnitude, \(\delta/(\pi^2/2) \approx 3.3 \times 10^{-4}\), falling
  within the natural order-of-magnitude range for discrete corrections
  of hypercube lattices (\(\sim 10^{-3}\)).
\end{itemize}

\subsubsection{6.3 Limited Claims}\label{limited-claims}

This document claims:

\begin{itemize}
\tightlist
\item
  The \textbf{sign and magnitude} of \(\delta\) are not inconsistent
  with the geometric structures of §3--§4.
\end{itemize}

This document does NOT claim:

\begin{itemize}
\tightlist
\item
  That the specific value \(\delta = 0.001619\) has been derived from
  geometry.
\item
  A complete numerical derivation.
\item
  That \(\delta\) cannot be explained under interpretive frameworks
  other than this one.
\end{itemize}

\begin{center}\rule{0.5\linewidth}{0.5pt}\end{center}

\subsection{§7 Limitations and Open Questions of This
Document}\label{limitations-and-open-questions-of-this-document}

\subsubsection{7.1 Unformalized Mappings}\label{unformalized-mappings}

A clear gap exists in this document:

\textbf{Gap G1}: A normalization rule converting the \(R^2\) deficit
term into the coefficient term \(\delta\) of the α equation is not
formulated.

The second term in (3.4) is an \(R^2\)-order volume correction, while
\(\delta\) in §6 is a dimensionless coefficient in the α equation. To
connect the two, a mapping rule is needed between:

\begin{itemize}
\tightlist
\item
  Volume corrections at the specific scale \(R = 3\) (where Paper 7
  yields 137), and
\item
  The self-consistency condition of the α equation.
\end{itemize}

This document does not provide this rule.

\subsubsection{7.2 Higher-Order
Corrections}\label{higher-order-corrections}

The third and subsequent terms (\(O(R)\), \(O(1)\)) of (3.4) are not
computed in this document. The complete asymptotic expansion takes the
form:

\[G(R) = \frac{16\pi}{3}R^3 - \left(\frac{3\pi^2}{4} + \frac{9\pi}{2}\right)R^2 + c_1 R + c_0 + o(1),\]

where \(c_1, c_0\) can also be analytically computed; however,
evaluation of these contributions is beyond the scope of this document.

\subsubsection{7.3 Possibility of Other
Interpretations}\label{possibility-of-other-interpretations}

This document presents only \textbf{one geometric interpretation}. The
possibility that the same numerical value emerges naturally from other
geometric structures (e.g., 4D cross-polytope packing, 4D root system
lattices, etc.) is not excluded.

In particular, in the framework of structural correspondence with Wilson
standard lattice gauge theory shown in Paper 8 {[}BH8{]}, similar
correction terms may be obtained via different routes.

\begin{center}\rule{0.5\linewidth}{0.5pt}\end{center}

\subsection{§8 Implications for Paper 7}\label{implications-for-paper-7}

\subsubsection{8.1 Positioning as
Reinforcement}\label{positioning-as-reinforcement}

The contribution of this document is limited. This document does NOT:

\begin{itemize}
\tightlist
\item
  \textbf{Modify} the numerical fact of 8.7 ppb precision in Paper 7
\item
  \textbf{Justify} the choice of \(\pi^2/2\) as a coefficient (it merely
  presents an alternative reading of the same numerical value)
\item
  \textbf{Solve} the residual problem (no complete derivation of
  \(\delta\) is provided)
\end{itemize}

However, this document does provide:

\begin{itemize}
\tightlist
\item
  A geometric viewpoint that \(\pi^2/2\) admits reading as a
  scale-invariant boundary measure core
\item
  The observation that, under this viewpoint, the sign and magnitude of
  the residual are interpretable
\item
  A partial response to the concerns about Paper 7 (Concerns A and B in
  §1.3)
\end{itemize}

\subsubsection{8.2 Effect on Defense
Lines}\label{effect-on-defense-lines}

This document does \textbf{NOT depend on} the claims of Paper 7:

\begin{itemize}
\tightlist
\item
  The 8.7 ppb observation of Paper 7 does not change regardless of
  acceptance or rejection of this document's interpretation.
\item
  Even if this document is partially rejected, Paper 7 remains
  unaffected.
\item
  This document is a reinforcement attempt and does not modify the
  logical foundation of Paper 7.
\end{itemize}

\subsubsection{8.3 Problems Posed for Subsequent
Research}\label{problems-posed-for-subsequent-research}

Problems posed by this document:

\begin{enumerate}
\def\labelenumi{\arabic{enumi}.}
\tightlist
\item
  Resolution of Gap G1 (§7.1): formulation of the rule converting the
  \(R^2\) deficit term to the α equation
\item
  Analytical computation of higher-order terms (\(c_1, c_0\))
\item
  Derivation of an isomorphic correction term from other geometric
  structures (4D cross-polytope, root system lattice, etc.)
\item
  Re-interpretation in the framework of Wilson structural correspondence
  {[}BH8{]}
\end{enumerate}

\begin{center}\rule{0.5\linewidth}{0.5pt}\end{center}

\subsection{§9 Conclusion}\label{conclusion}

This document provided one geometric interpretation of \(\pi^2/2\) in
the α identity \(\alpha^{-1} = 137 + (\pi^2/2)\alpha\) of Paper 7
{[}BH7{]}:

\begin{quote}
\textbf{Reading \(\pi^2/2\) as the normalization core
\(A_{S^3}(R)/(4R^3)\) of the 4D ball boundary \(S^3\) measure allows the
structure of Paper 7 to admit a scale-invariant interpretation.}
\end{quote}

Under this interpretation, the volume deficit \(G(R)\) of the
cube-covering of the 4D ball admits the expansion:

\[G(R) = \frac{16\pi}{3}R^3 - \left(\frac{3\pi^2}{4} + \frac{9\pi}{2}\right)R^2 + O(R),\]

where the second coefficient is necessarily negative under our
approximation model. This is consistent in sign and magnitude with the
back-calculated \(\delta \approx 0.001619\) from observation.

However, this document:

\begin{itemize}
\tightlist
\item
  Is an observation paper, not a proof paper
\item
  Does not provide a complete derivation of \(\delta\)
\item
  Does not exclude other interpretations
\item
  Does not modify the claims of Paper 7
\end{itemize}

The contribution of this document is limited to: presenting one
geometric viewpoint on the structure of Paper 7, and observing that this
viewpoint is not numerically inconsistent.

When future researchers, perhaps a hundred years from now, provide a
more complete geometric explanation of the structure of Paper 7, this
document may be referenced as an intermediate observation. That is
sufficient.

\begin{center}\rule{0.5\linewidth}{0.5pt}\end{center}

\subsection{Appendix A: Numerical Experiment
Implementation}\label{appendix-a-numerical-experiment-implementation}

\begin{Shaded}
\begin{Highlighting}[]
\ImportTok{import}\NormalTok{ numpy }\ImportTok{as}\NormalTok{ np}

\KeywordTok{def}\NormalTok{ count\_full\_cells\_4d(R):}
    \CommentTok{"""Count unit hypercube cells fully contained in 4D ball of radius R"""}
\NormalTok{    R\_sq }\OperatorTok{=}\NormalTok{ R }\OperatorTok{*}\NormalTok{ R}
\NormalTok{    count }\OperatorTok{=} \DecValTok{0}
\NormalTok{    R\_int }\OperatorTok{=} \BuiltInTok{int}\NormalTok{(R) }\OperatorTok{+} \DecValTok{1}
    \ControlFlowTok{for}\NormalTok{ n1 }\KeywordTok{in} \BuiltInTok{range}\NormalTok{(}\OperatorTok{{-}}\NormalTok{R\_int, R\_int }\OperatorTok{+} \DecValTok{1}\NormalTok{):}
        \ControlFlowTok{for}\NormalTok{ n2 }\KeywordTok{in} \BuiltInTok{range}\NormalTok{(}\OperatorTok{{-}}\NormalTok{R\_int, R\_int }\OperatorTok{+} \DecValTok{1}\NormalTok{):}
            \ControlFlowTok{for}\NormalTok{ n3 }\KeywordTok{in} \BuiltInTok{range}\NormalTok{(}\OperatorTok{{-}}\NormalTok{R\_int, R\_int }\OperatorTok{+} \DecValTok{1}\NormalTok{):}
                \ControlFlowTok{for}\NormalTok{ n4 }\KeywordTok{in} \BuiltInTok{range}\NormalTok{(}\OperatorTok{{-}}\NormalTok{R\_int, R\_int }\OperatorTok{+} \DecValTok{1}\NormalTok{):}
                    \CommentTok{\# Check if the farthest of the 8 vertices is inside the ball}
\NormalTok{                    max\_d\_sq }\OperatorTok{=} \DecValTok{0}
                    \ControlFlowTok{for}\NormalTok{ s1 }\KeywordTok{in}\NormalTok{ [}\OperatorTok{{-}}\FloatTok{0.5}\NormalTok{, }\FloatTok{0.5}\NormalTok{]:}
                        \ControlFlowTok{for}\NormalTok{ s2 }\KeywordTok{in}\NormalTok{ [}\OperatorTok{{-}}\FloatTok{0.5}\NormalTok{, }\FloatTok{0.5}\NormalTok{]:}
                            \ControlFlowTok{for}\NormalTok{ s3 }\KeywordTok{in}\NormalTok{ [}\OperatorTok{{-}}\FloatTok{0.5}\NormalTok{, }\FloatTok{0.5}\NormalTok{]:}
                                \ControlFlowTok{for}\NormalTok{ s4 }\KeywordTok{in}\NormalTok{ [}\OperatorTok{{-}}\FloatTok{0.5}\NormalTok{, }\FloatTok{0.5}\NormalTok{]:}
\NormalTok{                                    d\_sq }\OperatorTok{=}\NormalTok{ (n1 }\OperatorTok{+}\NormalTok{ s1)}\OperatorTok{**}\DecValTok{2} \OperatorTok{+}\NormalTok{ (n2 }\OperatorTok{+}\NormalTok{ s2)}\OperatorTok{**}\DecValTok{2} \OperatorTok{+} \OperatorTok{\textbackslash{}}
\NormalTok{                                           (n3 }\OperatorTok{+}\NormalTok{ s3)}\OperatorTok{**}\DecValTok{2} \OperatorTok{+}\NormalTok{ (n4 }\OperatorTok{+}\NormalTok{ s4)}\OperatorTok{**}\DecValTok{2}
                                    \ControlFlowTok{if}\NormalTok{ d\_sq }\OperatorTok{\textgreater{}}\NormalTok{ max\_d\_sq:}
\NormalTok{                                        max\_d\_sq }\OperatorTok{=}\NormalTok{ d\_sq}
                    \ControlFlowTok{if}\NormalTok{ max\_d\_sq }\OperatorTok{\textless{}=}\NormalTok{ R\_sq:}
\NormalTok{                        count }\OperatorTok{+=} \DecValTok{1}
    \ControlFlowTok{return}\NormalTok{ count}

\KeywordTok{def}\NormalTok{ G\_R(R):}
    \CommentTok{"""Compute volume deficit G(R)"""}
\NormalTok{    V\_4 }\OperatorTok{=}\NormalTok{ (np.pi}\OperatorTok{**}\DecValTok{2} \OperatorTok{/} \DecValTok{2}\NormalTok{) }\OperatorTok{*}\NormalTok{ R}\OperatorTok{**}\DecValTok{4}
\NormalTok{    N\_full }\OperatorTok{=}\NormalTok{ count\_full\_cells\_4d(R)}
    \ControlFlowTok{return}\NormalTok{ V\_4 }\OperatorTok{{-}}\NormalTok{ N\_full}

\CommentTok{\# Example}
\ControlFlowTok{for}\NormalTok{ R }\KeywordTok{in}\NormalTok{ [}\DecValTok{50}\NormalTok{, }\DecValTok{100}\NormalTok{, }\DecValTok{200}\NormalTok{, }\DecValTok{300}\NormalTok{]:}
\NormalTok{    g }\OperatorTok{=}\NormalTok{ G\_R(R)}
    \BuiltInTok{print}\NormalTok{(}\SpecialStringTok{f"R=}\SpecialCharTok{\{}\NormalTok{R}\SpecialCharTok{\}}\SpecialStringTok{: G(R)/R\^{}3 = }\SpecialCharTok{\{}\NormalTok{g}\OperatorTok{/}\NormalTok{R}\OperatorTok{**}\DecValTok{3}\SpecialCharTok{:.4f\}}\SpecialStringTok{, "}
          \SpecialStringTok{f"[G(R) {-} 16π/3·R³]/R² = }\SpecialCharTok{\{}\NormalTok{(g }\OperatorTok{{-}}\NormalTok{ (}\DecValTok{16}\OperatorTok{*}\NormalTok{np.pi}\OperatorTok{/}\DecValTok{3}\NormalTok{)}\OperatorTok{*}\NormalTok{R}\OperatorTok{**}\DecValTok{3}\NormalTok{)}\OperatorTok{/}\NormalTok{R}\OperatorTok{**}\DecValTok{2}\SpecialCharTok{:.2f\}}\SpecialStringTok{"}\NormalTok{)}
\end{Highlighting}
\end{Shaded}

Note: The above is conceptual implementation. For \(R = 300\), the total
number of loops becomes \(4 \cdot 600^4 \approx 5 \times 10^{11}\),
which is impractical. Implementation must optimize the search region
near the inner edge of the ball, or use Monte Carlo estimation.

\begin{center}\rule{0.5\linewidth}{0.5pt}\end{center}

\subsection{Appendix B: Details of Sphere
Integrals}\label{appendix-b-details-of-sphere-integrals}

The detailed derivations of the sphere integrals used in Lemmas 3.1 and
3.2 are completed in the main text. Recap:

\subsubsection{\texorpdfstring{B.1
\(\int_{S^3} |u_1| \, d\Omega = 8\pi/3\)}{B.1 \textbackslash int\_\{S\^{}3\} \textbar u\_1\textbar{} \textbackslash, d\textbackslash Omega = 8\textbackslash pi/3}}\label{b.1-int_s3-u_1-domega-8pi3}

The volume element of \(S^3\) can be parameterized as
\(u_1 = \cos\theta_1\) with
\(\sin^2\theta_1 \cdot d\theta_1 \cdot 4\pi\) (surface area of \(S^2\)).

\[\int_{S^3} |u_1| \, d\Omega = 4\pi \int_0^\pi |\cos\theta| \sin^2\theta \, d\theta = 8\pi \int_0^{\pi/2} \cos\theta \sin^2\theta \, d\theta = \frac{8\pi}{3}.\]

\subsubsection{\texorpdfstring{B.2
\(\int_{S^3} |u_1 u_2| \, d\Omega = \pi\)}{B.2 \textbackslash int\_\{S\^{}3\} \textbar u\_1 u\_2\textbar{} \textbackslash, d\textbackslash Omega = \textbackslash pi}}\label{b.2-int_s3-u_1-u_2-domega-pi}

With \(u_1 = \cos\theta_1\), \(u_2 = \sin\theta_1 \cos\theta_2\):

\[\int |u_1 u_2| \, d\Omega = 2\pi \int_0^\pi |\cos\theta_1| \sin^3\theta_1 \, d\theta_1 \cdot \int_0^\pi |\cos\theta_2| \sin\theta_2 \, d\theta_2 = 2\pi \cdot \frac{1}{2} \cdot 1 = \pi.\]

\subsubsection{\texorpdfstring{B.3
\(\int_{S^3} u_i^2 \, d\Omega = \pi^2/2\)}{B.3 \textbackslash int\_\{S\^{}3\} u\_i\^{}2 \textbackslash, d\textbackslash Omega = \textbackslash pi\^{}2/2}}\label{b.3-int_s3-u_i2-domega-pi22}

By symmetry,
\(\int u_i^2 = (1/4) \int \sum_j u_j^2 = (1/4) \cdot \text{Vol}(S^3) = (1/4) \cdot 2\pi^2 = \pi^2/2\).

\begin{center}\rule{0.5\linewidth}{0.5pt}\end{center}

\subsection{References}\label{references}

\begin{itemize}
\tightlist
\item
  {[}BH7{]} Kihara, N., \emph{A Geometric Identity for the
  Fine-Structure Constant: From the 4D Unit Ball Volume and its
  Cube-Packing Deficit}, Concept DOI:
  \href{https://doi.org/10.5281/zenodo.19869266}{10.5281/zenodo.19869266}
  (2026).
\item
  {[}BH8{]} Kihara, N., \emph{Chain Complex Structure on the 4D
  Hypercubic Lattice: Structural Correspondence between Kihara
  Cube-Packing and Wilson Lattice Gauge Theory}, Concept DOI:
  \href{https://doi.org/10.5281/zenodo.19880467}{10.5281/zenodo.19880467}
  (2026).
\item
  {[}Wilson74{]} Wilson, K. G., \emph{Confinement of quarks}, Phys.
  Rev.~D 10, 2445 (1974).
\item
  {[}CODATA22{]} CODATA Recommended Values of the Fundamental Physical
  Constants: 2022.
\end{itemize}

\begin{center}\rule{0.5\linewidth}{0.5pt}\end{center}

Author: Noriaki Kihara WF System Co., Ltd.~/ ORCID:
\href{https://orcid.org/0009-0004-6753-4020}{0009-0004-6753-4020} Paper
DOI (to be assigned): not yet acquired License: CC BY 4.0

\end{document}
